\clearpage
\fancyhead[L]{\small PASQUICH. STERNWARTE IN OFEN.}
\hypertarget{native-502-lower}{}\label{native-502-lower}
\begin{center}
\small Göttingische gelehrte Anzeigen. Stück 85. Seite 841..843. 1808 Mai 28.\\[3pt]
\rule{.72\linewidth}{.3pt}
\end{center}

Unter dem Titel: \emph{Rechenschaft von meinen Vorschlägen zur Beförderung der Astronomie auf der königl. Universitäts-Sternwarte in Ofen}, hat der Director dieser Sternwarte, Hr. \textsc{Joh. Pasquich}, eine kleine Schrift drucken lassen, die über das zu hoffende künftige Aufblühen der Astronomie in Ungarn sehr erfreuliche Nachrichten enthält.

{\rightskip=0pt plus 3em Die seit 1780 in Ofen befindliche Sternwarte ist weder in Rücksicht der Bauart, noch der Instrumente, die sie besitzt, dem heutigen Zustande der Astronomie angemessen: dies war um so mehr zu bedauern, da den Statuten der Universität nach für die Sternwarte ausser dem Director noch zwei Adjuncten und ein Wärter unterhalten werden müssen. Die königliche Ungarische Statthalterei, an deren Spitze der erleuchtete Erzherzog Palatinus steht, sah ein, wie wichtig dieses Etablissement für die Wissenschaft und das Land werden könne, wenn hier thätige und ganz für ihre Wissenschaft lebende Männer mit bessern Hülfsmitteln ausgerüstet würden. Hr. Prof. \textsc{Pasquich} erhielt daher gleich nach seiner Anstellung als Astronom bei dieser Sternwarte den höchsten Befehl, Vorschläge zur Anschaffung neuer Instrumente zu thun: er that dies nach dem Grundsatze, dass es hier wirklich auf echte Beförderung der Wissenschaft abgesehen sei, zu welcher alle cultivirten Nationen Europas beizutragen sich verpflichtet halten. Seine Vorschläge wurden auch sowohl von dem Erzherzog Palatinus und der königl Statthalterei, als von dem Kaiser selbst, sogleich ohne alle Einschränkung genehmigt, und die ungesäumte Bestellung der Instrumente befohlen. Er wandte sich deshalb an den Artillerie-Hauptmann \textsc{Reichenbach} in München, der bekanntlich seit einigen Jahren dort eine Werkstätte zur Verfertigung mathematischer Instrumente errichtet hat, und nach mehreren Proben den ersten englischen Künstlern den Rang streitig machen zu können scheint. Es wurden bestellt: ein dreifussiger Repetitionskreis mit dreissigzolligem Azimutalkreise und silbernem Limbus, ein sechsfussiges Mittagsfernrohr, eine astronomische Secunden-Pendeluhr, eine astronomische Reise-Halbesecunden-Pendeluhr, ein achtzehnzolliger 
\SourceLeaf{503}
 astronomischer Kreis, ein zwölfzolliger Kreis zu terrestrischen Messungen, ein achtfussiges Fernrohr, ein Aequatoreal; welches alles, einige Reparaturen an andern Instrumenten noch mit inbegriffen, zu dem mässigen Preise von 7210 kaiserl. Gulden oder 8652 Gulden Reichswährung bedungen wurde. Ausserdem wurde noch eine astronomische Secunden-Pendeluhr bei dem Bergrath \textsc{Seyffert} in Dresden zu 360 Thaler bestellt. Die Ausführung jener Instrumente wurde durch den Krieg zwar unterbrochen, ist aber jetzt bis auf das achtfussige Fernrohr und das Aequatoreal ganz vollendet, und ein Theil der Instrumente bereits in Ofen angelangt, die übrigen werden täglich erwartet.\par}

Es bleibt daher jetzt nichts zu wünschen übrig, als dass dieser schöne Apparat bald ein würdiges Local finden möge, um zum Besten der Wissenschaft und zur Ehre der Ungrischen Nation und der hohen Beförderer angewandt werden zu können. Hr. \textsc{Pasquich} hat bereits einen schicklichen Platz für die neue Sternwarte vorgeschlagen, und nur einige besondre Umstände machen es noch ungewiss, ob man sich für denselben entscheiden wird. Auf alle Fälle darf man sicher erwarten, dass in kurzem Ungarn mit einem Tempel der Urania geschmückt sein wird, wozu man der Ungrischen Nation und der Wissenschaft selbst wird Glück wünschen können.
